\documentclass[9pt,twocolumn]{extarticle}
\usepackage{journaloflove-sci}
% ============================================================
% EDIT THIS FILE ONLY / 只改这里即可同步更新两份学报
% ============================================================
\newcommand{\PartnerA}{姓名 A}
\newcommand{\PartnerB}{姓名 B}
\newcommand{\FirstMeetingDate}{20XX 年 XX 月 XX 日}
\newcommand{\RelationshipDate}{20XX 年 XX 月 XX 日}
\newcommand{\AnnouncementDate}{20XX 年 XX 月 XX 日}
\newcommand{\WeddingDate}{20XX 年 XX 月 XX 日}
\newcommand{\WeddingCity}{城市}
\newcommand{\WeddingVenue}{婚礼场地 / 酒店 / 草坪}
\newcommand{\ContactEmail}{xxx@foreverlove.com}
\newcommand{\OurMotto}{此后，日常即共同研究；岁月为长期追踪。}
\newcommand{\DaysObserved}{XXXX}
\newcommand{\SharedTrips}{XX}
\newcommand{\SharedMeals}{XXXX+}
\newcommand{\FutureHorizon}{终身}


\setRunningAuthor{\PartnerA{} and \PartnerB{}}
\setJOLMeta{Journal of Love 1314 (2026) 520--521}
\setJOLCoverSub{WEDDING ISSUE}
\setJOLCoverIssue{VOL. 1314 / NO. 520}
\setJOLTitle{From First Encounter to a Shared Life: A Longitudinal Study of Pair Bonding and Marriage Formalization}
\setJOLAuthors{\PartnerA{}\textsuperscript{a}, \PartnerB{}\textsuperscript{a,*}}
\setJOLAffilA{Joint Laboratory of Everyday Life, \WeddingCity{}, Earth}
\setJOLKeywords{Marriage\\Pair bonding\\Longitudinal relationship\\Wedding\\Shared life}
\setJOLAbstract{This study documents the longitudinal formation and formalization of a two-person partnership between \PartnerA{} and \PartnerB{}. The observation window begins on \FirstMeetingDate{} and includes repeated high-frequency interactions, shared experiences and joint decisions. The evidence indicates persistent reciprocal preference, increasing coordination and a durable intention to share future life. The wedding scheduled for \WeddingDate{} at \WeddingVenue{} is interpreted as the formal institutionalization of this long-term collaboration rather than its endpoint. The follow-up horizon is defined as \FutureHorizon{}.}
\setJOLCorrespondence{Joint Laboratory of Everyday Life, \WeddingCity{}, Earth.}
\setJOLEmail{\ContactEmail}
\setJOLDOI{https://doi.org/10.1314/jlove.520.521}
\setJOLReceived{Received \FirstMeetingDate{}; Received in revised form \RelationshipDate{}; Accepted \WeddingDate{}}
\setJOLOnline{Available online \WeddingDate{}}
\setJOLCopyright{1314-0520/© 2026 Cupid Press. All hearts reserved.}
\setJOLFootnoteText{The longitudinal dataset consists of private messages, photographs, shared itineraries and everyday memories. Selected observations may be disclosed only with mutual consent.}

\begin{document}
\setcounter{page}{520}
\makeJOLfrontmatter

\section{Introduction}
长期伴侣关系可以被理解为一条逐步形成的共同轨迹。对本文的两个观察对象而言，\FirstMeetingDate{} 构成初始时点。此后，两条原本相互独立的生活路径开始不断出现交集：共同用餐、旅行、讨论、争执后的修复，以及大量无法被结构化数据完整编码的日常细节。

本文借用学术论文的形式记录一项已经得到双方共同确认的事实：在长期互动与持续更新之后，两位作者决定把彼此纳入未来的基准情景，并于 \WeddingDate{} 在 \WeddingVenue{} 举行婚礼。\textsuperscript{1}

\section{Materials and methods}
\subsection{Sample and observation window}
研究样本包含两个核心个体，记为 $i\in\{A,B\}$。观察期自 \FirstMeetingDate{} 起，覆盖约 \DaysObserved{} 个自然日。主要信息来源包括共同经历、即时沟通、线下陪伴、旅行记录以及双方对重大生活决策的联合讨论。

\subsection{Conceptual specification}
为保持论文语体，我们将长期关系状态写成一个极简动态过程：
\begin{equation}
R_{t+1}=R_t+\alpha T_t+\beta C_t+\gamma S_t-\delta D_t+\varepsilon_t,
\end{equation}
其中 $R_t$ 表示关系稳定度，$T_t$ 表示信任，$C_t$ 表示沟通质量，$S_t$ 表示共同支持，$D_t$ 表示需要共同解决的分歧。该式仅用于叙事展示，并不主张对爱情进行真正的结构估计。

\clearpage
\section{Results and discussion}
长期追踪显示，双方关系并未因观察期延长而衰减。随着共同经历增加，生活协同性与未来规划的一致性持续积累。若把长期共同生活意愿作为主要结局变量，则该变量已获得双向一致观测。

\begin{JOLtable}{Selected descriptive evidence from the relationship panel}
\begin{tabularx}{\columnwidth}{@{}l>{\raggedleft\arraybackslash}X@{}}\toprule
Indicator & Observation\\\midrule
Observation days & \DaysObserved{}\\
Shared trips & \SharedTrips{}\\
Shared meals & \SharedMeals{}\\
Mutual commitment & Confirmed\\
Wedding decision & Jointly approved\\
Follow-up horizon & \FutureHorizon{}\\\bottomrule
\end{tabularx}
\end{JOLtable}

关系中的高频细节并不总适合量化。很多真正具有解释力的变量——被理解的感觉、困难时期的陪伴、无需说明的默契——往往不能被准确编码。因此，本文的“论文式”表达是一种仪式化叙事：借用研究设计的严谨外壳，保存一段私人关系的公开版本。

\JOLphoto{assets/wedding-photo.jpg}{Wedding photograph.}

\subsection{Wedding as an event}
婚礼安排如下：\WeddingDate{}，地点为 \WeddingCity{} 的 \WeddingVenue{}。在本文框架下，这一事件被视作一项公开、正式且具有社会见证的制度化安排。它不是长期关系的终点，而是后续观察期的起点。

\begin{JOLtable}{Wedding event information}
\begin{tabularx}{\columnwidth}{@{}lX@{}}\toprule
Field & Information\\\midrule
Couple & \PartnerA{} $\times$ \PartnerB{}\\
Date & \WeddingDate{}\\
City & \WeddingCity{}\\
Venue & \WeddingVenue{}\\
Contact & \ContactEmail{}\\\bottomrule
\end{tabularx}
\end{JOLtable}

\section{Conclusion}
本文记录了一项简单但重要的结果：\PartnerA{} 与 \PartnerB{} 已决定结婚。所有过去的数据用于理解相遇，所有未来的数据将用于共同生活。后续研究窗口为 \FutureHorizon{}。

\section*{CRediT authorship contribution statement}
\PartnerA{}: Conceptualization, Investigation, Daily-life data collection, Emotional support, Writing -- original draft. \PartnerB{}: Conceptualization, Investigation, Daily-life data collection, Emotional support, Writing -- review \& editing. Contribution is declared equal.

\section*{Declaration of competing interest}
The authors declare no competing interests, except for a persistent bilateral preference for one another.

\section*{Data availability}
The underlying dataset consists of private messages, photographs, shared itineraries and everyday memories and is not publicly available for privacy reasons.

\section*{Acknowledgements}
感谢双方家人、朋友以及所有见证这段关系的人。婚礼当日的每一次到场，都将被视为对本研究最温柔的同行评议。

\begin{thebibliography}{9}\fontsize{6.7}{8.2}\selectfont
\bibitem{memory} \PartnerA{} and \PartnerB{}, Shared Memory Dataset: Private longitudinal records, unpublished private archive, \FirstMeetingDate{}--present.
\bibitem{future} \PartnerA{} and \PartnerB{}, Protocol for lifelong follow-up, jointly approved research plan, \WeddingDate{}.
\end{thebibliography}
\end{document}
